\chapter{Implications for OpenFlight}
\label{ch:implications}

OpenFlight's current architecture --- OPS243-A 24\,GHz CW Doppler with
rolling-buffer I/Q capture, a hardware sound trigger, and two K-LD7
interferometric angle radars, on a Raspberry Pi 5 --- is a minimal but
honest instance of the commercial radar architecture. This chapter
distills the review into design guidance, ordered by leverage.

\section{Ball data}

\begin{enumerate}
\item \textbf{Ball speed} is already the strongest channel: the Doppler
line is unambiguous, and the OPS243-A's $\pm0.5\%$ class accuracy matches
commercial practice. Add the aspect-angle (cosine) correction using the
measured launch angles --- at 3--5\,ft behind the tee with a driver launch
of $12\degs$ the radial underestimate is small but systematic.
\item \textbf{Launch angles}: treat the K-LD7 bursts as a short
trajectory, not a single detection --- an EKF smoother over the ring
buffer (\cref{sec:ekf}) with back-extrapolation to the impact timestamp
is the single highest-leverage software upgrade for angle quality.
\item \textbf{Spin rate}: the current $\sim$50--60\% detection rate from
I/Q sideband analysis is consistent with the physics: sideband SNR
depends on ball-surface asymmetry and observation time
(\cref{sec:radarspin}). Improvements in order of cost: longer
observation windows (later trigger cutoff); cepstrum/harmonic-product
comb estimation with the patent-style equal-spacing qualification;
logo-ball or marked-ball guidance for indoor use; and honest fallback ---
report estimated spin (from club type + speed + loft priors, clearly
flagged) when no comb is found, as Garmin does. \emph{Patent caution:}
the harmonic-sideband method is claimed by US8845442 until $\sim$2029
(US member); for a hobby project this is a risk-management judgment, but
the roadmap should note the family's expiry dates
(\cref{sec:fto}).
\item \textbf{Spin axis} is currently out of reach of the radar: indoor
flights are too short for trajectory inversion, and the direct
receiver-pair method is FlightScope-patented to $\sim$2035. The practical
route is optical (below) or a flagged model estimate from face-to-path.
\end{enumerate}

\section{Club data}

\begin{enumerate}
\item \textbf{Club speed}: define and document the reference point. The
OPS243-A sees a smear of head/shaft returns; gating the pre-impact
spectrum and taking a fixed percentile of the velocity distribution
(rather than the maximum) approximates a stable reference and avoids the
toe-speed inflation that plagues naive processors
(\cref{sec:radarclub}). Validate the smash factor against the
loft-appropriate ceiling (\cref{sec:smash}) and flag violations instead
of displaying them.
\item \textbf{Club path} (K-LD7 horizontal) and \textbf{attack angle}
(vertical, if geometry allows) are legitimate direct measurements ---
the same primitives the Garmin R10 measures. Alignment calibration
dominates their accuracy.
\item \textbf{Face angle}: implement the D-plane inversion
(\cref{eq:faceinversion}) with the loft-dependent weight
(\cref{eq:obliqueness}) --- this is exactly what radar-only commercial
units report --- but label it \emph{derived}, propagate its error
($1.15\times$ launch-direction bias), and suppress it on detected
off-center strikes once impact location is available. The D-plane
forward model doubles as the simulator's launch generator, so one
well-tested module serves both directions.
\item \textbf{Impact location} requires optics. The expired Acushnet
stereo art (\cref{sec:patentacushnet}) and Wintriss mono-camera art
(\cref{sec:wintriss}) provide complete public-domain recipes; a
PiTrac-style Pi Global Shutter camera with IR strobing is the natural
hardware. Even a single overhead camera reading face stickers
(Acushnet US6758759 recipe, expired) would convert OpenFlight's face
angle from derived to measured.
\end{enumerate}

\section{Architecture roadmap}

\begin{enumerate}
\item \textbf{Phase 1 (radar-only, current)}: EKF smoothing, comb-based
spin with honest fallback, D-plane-derived face data with flags,
smash-factor sanity gates, alignment-calibration procedure. Ships the
Garmin-R10 feature set with open data.
\item \textbf{Phase 2 (optical spin/impact module)}: one or two Pi
Global Shutter cameras + IR strobe PCB (PiTrac-proven, $\sim$\$100
incremental). Camera measures launch angles, 3D spin (dimple
registration, public-domain per expired US7292711), and impact location;
radar keeps trigger, speed, and outdoor robustness. This is the same
fusion every commercial vendor converged on --- but built from the
expired-art side of the patent map. \emph{Care:} avoid the still-active
claims around integrated ball-find/LED-guidance UX (to $\sim$2027) and
radar-steered image search windows (FlightScope US10338209).
\item \textbf{Phase 3 (fusion)}: a single EKF consuming OPS243-A speed,
K-LD7 angles, and camera fixes, with per-sensor covariances; flight
model per \cref{ch:flight} with versioned coefficients; validation
protocol against the MLM2PRO per \cref{ch:accuracy}.
\item \textbf{Phase 4 (measured club delivery)}: IWR6843 with custom
chirp firmware feeding the screw-theoretic rigid-body estimator of
\cref{app:screw} --- per-detection Doppler rows solved for the club's
twist, smoothed on $SE(3)$, and projected into club speed at a declared
reference point, measured path and attack angle, closure rate, and an
ISA-defined swing plane. This moves path and attack angle from
\emph{derived} to \emph{measured} in \cref{tab:hierarchy}; face angle
and impact location remain with the optical module.
\end{enumerate}

\section{Reporting honesty as a feature}

The single clearest lesson of the review: commercial marketing blurs the
measured/derived boundary, and the validation literature repeatedly
punishes the derived quantities. An open-source instrument can win trust
by doing the opposite --- tagging every reported parameter with its
provenance (measured / derived / estimated), its uncertainty, and the
model version used. \Cref{tab:hierarchy} is effectively the schema for
that tagging.
